\chapter{Medical and surgical abortions}

\section*{Definition and Overview}
Medical and surgical abortions are the two primary clinical modalities utilised to terminate an unwanted, unintended, or medically compromised pregnancy. Termination of pregnancy is a safe, highly regulated, and essential component of comprehensive reproductive healthcare. A medical abortion involves the sequential administration of oral or vaginal medications (typically mifepristone and misoprostol) to induce uterine contractions and expel the pregnancy tissue, commonly restricted to the first trimester. A surgical abortion is a minor gynaecological procedure, such as suction aspiration (vacuum curettage) or dilation and evacuation (D\&E), to mechanically empty the uterine cavity. The selection of the abortion method depends on gestational age, maternal medical history, contraindications, personal preference, and local healthcare accessibility.

\section*{Epidemiology}
Abortion is one of the most common medical procedures performed globally, with millions of procedures carried out annually. In Australia, it is estimated that approximately 80,000 abortions are performed each year. The vast majority of terminations (over 90\%) occur during the first trimester (under 12 weeks of gestation). The introduction and registration of medical abortion (using mifepristone and misoprostol) has significantly transformed clinical practice, with medical abortions now accounting for a substantial and growing proportion of early terminations, particularly in primary care and regional settings. Surgical abortion remains the primary method for terminations in the second trimester and for patients who prefer a rapid, single-visit procedure with high efficacy.

\section*{Aetiology and Pathophysiology}
Abortions are performed for various social, personal, maternal, or fetal indications. The physiological mechanisms of the two methods are distinct:
\begin{itemize}
    \item \textbf{Medical abortion mechanism:} Mifepristone acts as a competitive progesterone receptor antagonist, blocking the action of progesterone, which leads to endometrial degeneration, cervical softening, and decidual detachment. Misoprostol, a synthetic prostaglandin E1 analogue administered 24 to 48 hours later, binds to myometrial cells to induce uterine contractions and cervical dilatation, resulting in the expulsion of the gestational sac.
    \item \textbf{Surgical abortion mechanism:} Vacuum aspiration (manual or electric) involves dilating the cervix (often preceded by cervical priming with misoprostol or osmotic dilators) and inserting a flexible or rigid cannula into the uterine cavity. Negative pressure is applied to safely evacuate the conceptus. For later gestations (typically past 14 weeks), dilation and evacuation (D\&E) utilises surgical instruments alongside suction to ensure complete removal of fetal and placental tissues.
\end{itemize}

\section*{Clinical Features}
Patients undergoing abortion experience symptoms related to the termination process, which vary between the two methods:
\begin{itemize}
    \item \textbf{Medical abortion features:} Moderate-to-severe lower abdominal cramping and pelvic pain, followed by moderate-to-heavy vaginal bleeding (often heavier than a normal period, with clots) as the pregnancy tissue is expelled. This is commonly accompanied by transient nausea, vomiting, diarrhoea, chills, or low-grade fever due to the systemic effects of misoprostol.
    \item \textbf{Surgical abortion features:} Mild lower abdominal cramping and light vaginal bleeding or spotting following the procedure, which typically resolve within 1 to 2 weeks. Postoperative pain is generally minimal and managed well with standard analgesia, with the procedure itself performed under local anaesthetic block, conscious sedation, or general anaesthesia.
    \item \textbf{Anaemic symptoms:} Fatigue, dizziness, or pallor if bleeding is exceptionally heavy or prolonged, prompting investigation of haemoglobin levels.
\end{itemize}

\section*{Differential Diagnosis}
During the post-abortion recovery phase, clinicians must distinguish normal post-abortion changes from potential complications. The differential diagnosis includes:
\begin{itemize}
    \item \textbf{Retained products of conception (RPOC) vs. Normal recovery:} RPOC presents with persistent heavy vaginal bleeding, escalating pelvic pain, and an enlarged, tender uterus, whereas normal recovery involves gradually tapering bleeding.
    \item \textbf{Haematometra (clot retention) vs. Uterine atony:} Haematometra is the accumulation of blood clots within the uterine cavity causing intense pain and a firm uterus, whereas uterine atony is a soft, poorly contracted uterus leading to heavy haemorrhage without significant mechanical obstruction.
    \item \textbf{Post-abortal pelvic infection vs. Prostaglandin-induced fever:} Pelvic infection (endometritis) presents with persistent fever (above 38 degrees Celsius), purulent or foul-smelling vaginal discharge, and uterine tenderness, whereas misoprostol causes a transient, self-limiting rise in temperature without pelvic tenderness.
    \item \textbf{Ectopic pregnancy vs. Incomplete abortion:} An ectopic pregnancy must be ruled out if no pregnancy tissue is identified during surgical aspiration or if beta-hCG levels do not fall adequately after a medical abortion.
\end{itemize}

\section*{When to Seek Medical Care}
Patients recovering from a medical or surgical abortion should follow up with their clinical team to confirm the completion of the procedure and discuss contraception.

\textbf{Call 000 (or local emergency services) for:}
\begin{itemize}
    \item \textbf{Haemorrhage:} Bleeding that is extremely heavy, defined as soaking through two or more maxi-sized sanitary pads per hour for two consecutive hours.
    \item \textbf{Severe, localised abdominal pain:} Intense, sharp, or worsening pelvic pain that is not relieved by analgesics, which may indicate uterine perforation, pelvic haematoma, or internal bleeding.
    \item \textbf{Anaphylaxis:} Swelling of the face, mouth, or throat, difficulty breathing, or severe skin rash.
    \item \textbf{Sepsis:} High fever, severe chills, rapid heart rate, or feeling extremely weak and confused.
    \item \textbf{Signs of shock:} Feeling faint, severe dizziness, cold and clammy skin, or loss of consciousness.
\end{itemize}

\section*{Diagnosis and Investigations}
Pre-procedure and post-procedure investigations ensure the safety and success of the abortion:
\begin{itemize}
    \item \textbf{Pelvic ultrasound:} The primary tool used pre-procedure to confirm intrauterine pregnancy and determine precise gestational age, and post-procedure if incomplete abortion or RPOC is suspected.
    \item \textbf{Serum beta-hCG level:} A baseline level is measured pre-procedure, and a follow-up test is performed 1 to 2 weeks after a medical abortion (or a high-sensitivity urine test at 3 to 4 weeks) to confirm a successful drop in hormone levels, indicating completion.
    \item \textbf{Maternal blood typing and Rh status:} Essential to identify Rh-negative patients, who require anti-D immunoglobulin within 72 hours of the abortion to prevent rhesus isoimmunisation in future pregnancies.
    \item \textbf{High vaginal swabs:} Screening for sexually transmitted infections (such as chlamydia and gonorrhoea) prior to surgical abortion to allow prophylactic antibiotic treatment and prevent ascending pelvic infection.
\end{itemize}

\section*{Management}
Clinical management involves providing the appropriate medications or performing the surgical procedure, combined with supportive care and contraception:
\begin{itemize}
    \item \textbf{Medical abortion regimen:} Administration of 200 mg oral mifepristone, followed 24 to 48 hours later by 800 mcg of misoprostol (buccally, sublingually, or vaginally), along with anti-emetics and strong analgesia (such as NSAIDs or codeine).
    \item \textbf{Surgical abortion procedure:} Performed in a day-surgery clinic under local anaesthetic (cervical block), sedation, or general anaesthesia. The cervix is dilated, vacuum suction is applied to evacuate the uterus, and the retrieved tissue is inspected to confirm completion.
    \item \textbf{Antibiotic prophylaxis:} Routine administration of antibiotics (such as azithromycin or doxycycline) around the time of surgical abortion to minimise infection risk.
    \item \textbf{Post-procedure pain relief:} Adequate pain management using ibuprofen or paracetamol, as well as heat packs applied to the lower abdomen.
    \item \textbf{Contraception initiation:} Immediate initiation of contraception post-abortion; an IUD can be inserted at the time of surgical abortion, or oral contraceptives/implants started on the day of medical or surgical termination.
\end{itemize}

\section*{Complications}
Both methods are highly safe, but complications can occur:
\begin{itemize}
    \item \textbf{Incomplete abortion / Retained products of conception:} Failure to completely expel the pregnancy, which occurs in 2\% to 5\% of medical abortions and less than 1\% of surgical abortions, requiring repeat medication or suction curettage.
    \item \textbf{Uterine perforation:} A rare complication of surgical abortion (less than 1 in 1,000 cases) where a surgical instrument accidentally passes through the uterine wall, potentially requiring laparoscopic assessment.
    \item \textbf{Cervical injury:} Tearing or laceration of the cervix during surgical dilatation, requiring sutures.
    \item \textbf{Infection:} Pelvic inflammatory disease or endometritis secondary to ascending vaginal bacteria, requiring antibiotics.
    \item \textbf{Uterine synechiae (Asherman syndrome):} Rare intrauterine scarring following surgical curettage, which can impact future fertility.
\end{itemize}

\section*{Prognosis and Outcomes}
The prognosis after both medical and surgical abortion is excellent. Over 95\% of medical abortions and 99\% of surgical abortions are successful without requiring further intervention. Physical recovery is rapid; normal activities can usually be resumed within 24 to 48 hours, and menstruation typically returns within 4 to 6 weeks. Future fertility is not adversely affected by uncomplicated medical or surgical abortions. Long-term psychological sequelae are rare, with the majority of patients experiencing a sense of relief, though access to post-abortion support services is beneficial for those who experience complex emotions or lack social support.

\section*{Prevention and Self-Care}
Post-abortal care is focused on preventing infections, promoting recovery, and establishing a future reproductive plan:
\begin{itemize}
    \item \textbf{Rest and recovery:} Limit strenuous physical activity and heavy lifting for 2 to 3 days post-procedure to minimise cramping and bleeding.
    \item \textbf{Avoiding infection:} Refrain from vaginal intercourse, swimming, bathing, and using tampons or menstrual cups for at least 1 to 2 weeks (or until bleeding stops) to protect the healing cervix.
    \item \textbf{Symptom monitoring:} Track the amount of bleeding and check temperature if feeling unwell.
    \item \textbf{Post-abortion follow-up:} Complete the scheduled follow-up (ultrasound, blood test, or specialised urine test) to ensure the pregnancy has been terminated.
    \item \textbf{Reliable contraception:} Adopt a highly effective contraceptive plan immediately after the procedure, as fertility can return as early as 10 days post-abortion.
\end{itemize}

\section*{Sources and Further Reading}
The following reputable organisations provide evidence-based information and guidelines on abortion:
\begin{itemize}
    \item MSI Reproductive Choices Australia. \textit{Information on abortion services, options, and support.} https://www.msichoices.org.au/
    \item Royal Australian and New Zealand College of Obstetricians and Gynaecologists (RANZCOG). \textit{Clinical guidelines and patient resources on termination of pregnancy.} https://ranzcog.edu.au/
    \item Healthdirect Australia. \textit{Overview of medical and surgical abortion pathways in Australia.} https://www.healthdirect.gov.au/
    \item NHS. \textit{Detailed guide to abortion methods, what to expect, and recovery.} https://www.nhs.uk/
    \item Better Health Channel. \textit{Reproductive health choices, abortion legalities, and support services.} https://www.betterhealth.vic.gov.au/
\end{itemize}
